\chapter*{Заключение}
\addcontentsline{toc}{chapter}{Заключение}

В ходе научно-исследовательской работы был разработан и экспериментально исследован модульный масштабируемый программный комплекс динамической навигации и мезоскопического агент-ориентированного моделирования транспортных потоков. 

Основные результаты работы:
\begin{enumerate}
    \item На основе анализа существующих подходов обоснована высокая вычислительная сложность классического автономного мультиагентного поиска путей на больших графах и предложена гибридная парадигма с декомпозицией функций на вычислительное ядро без сохранения состояния и внешний макроскопический модуль принятия решений. Разработана мезоскопическая модель симуляции, использующая пространственные очереди с физикой обратных волн затора на базе гидродинамической теории, которая гарантирует строгое выполнение физического принципа сохранения очередности движения (FIFO) через механизм линейной интерполяции временных задержек.
    
    \item Спроектирована архитектура навигационной системы с использованием механизма временных интервалов и разработано оптимизированное вычислительное ядро динамической маршрутизации на основе алгоритма поиска по ориентирам (ALT). Сжатое хранение топологии в формате CSR, топологическая пересортировка вдоль пространственной кривой Мортона для повышения локальности в кэше процессора и приоритетная очередь на основе векторизованной $8$-арной кучи ускорили поиск пути до $17{,}7$ раза по сравнению с алгоритмом Дейкстры. Тестирование на реальном дорожном графе Московской агломерации ($582$~тыс. узлов, $1~460$~тыс. рёбер) подтвердило пиковую пропускную способность в $5130$ запросов в секунду при $12$ вычислительных потоках.
    
    \item Разработан программный код агент-ориентированной системы, включающий асинхронный управляющий контур для триггерного перерасчета маршрутов агентов на основе динамических функций задержки (BPR) и фоновый модуль телеметрии. Использование безблокировочной очереди позволило изолировать сетевые задержки при передаче телеметрии от вычислительного ядра. В динамическом стресс-тесте управляющий контур показал высокую отказоустойчивость, успешно фильтруя избыточные запросы для сохранения нулевого пространственного лага агентов. Интегральное моделирование выявило неестественную пространственную концентрацию трафика (рост перегруженных сегментов до $3578$), что обосновывает необходимость доработки кинематики продвижения агентов.
    
    \item Проведено комплексное исследование масштабируемости разработанного комплекса. Сравнение различных очередей сообщений подтвердило преимущества безблокировочной архитектуры MPMC при межпоточном обмене. Оценка многопоточной масштабируемости выявила аппаратные ограничения пропускной способности оперативной памяти (memory wall) при интенсивном случайном поиске пути на многоядерных процессорах и доказала эффективность привязки потоков к ядрам процессора для повышения производительности на $13{,}6\%$.
\end{enumerate}

Разработанный программный комплекс координирует транспортные потоки уровня крупной агломерации на стандартных серверах среднего класса. Стек технологий на базе современного стандарта языка программирования C++26, концепций абстракций с нулевой стоимостью, сетевых сокетов и веб-технологий обеспечивает высокую производительность и удобную интеграцию системы в диспетчерские комплексы «Умного города».

В качестве перспективных направлений дальнейшего развития программного комплекса выделены две фундаментальные задачи. Во-первых, требуется фундаментальная доработка внутренней гидродинамической модели симулятора и кинематики движения агентов для устранения неестественной концентрации автомобилей. Во-вторых, предлагается преодоление теоретического ограничения на применение двунаправленного эвристического поиска ALT в сетях с динамическим бронированием емкостей. Внедрение гибридного предиктивного «Оракула» на базе глубокого обучения с подкреплением позволит аппроксимировать время прибытия на цель и инициализировать обратный поиск, что способно дополнительно повысить производительность маршрутизации на длинных дистанциях.

